\subsection{函子}

定义 3. —— 设 C 与 C' 为范畴。从 C 到 C' 的函子是箭图态射 \(\varphi\colon\mathrm{C}\to\mathrm{C}'\) ，满足 \(\varphi(fg)=\varphi(f)\varphi(g)\) （对 C 的每对可复合箭 \((f,g)\) ）。

设 C 与 C' 为范畴，\(\varphi\colon\mathrm{C}\to\mathrm{C}'\) 为函子。设 \(f\) 为 C 的箭，\(a\) 为其起点，\(b\) 为其靶；则 \(\varphi(f)\) 是 C' 的箭，起点为 \(\varphi(a)\) ，靶为 \(\varphi(b)\) 。

设 C、C'、C'' 为范畴，\(\varphi\colon\mathrm{C}\to\mathrm{C}'\) 与 \(\varphi'\colon\mathrm{C}'\to\mathrm{C}''\) 为函子。则 \(\varphi'\circ\varphi\) 是函子。

设 C 为范畴。则箭图态射 \(Id_{C}\) 是函子。

设 \(\varphi\colon\mathrm{C}\to\mathrm{C}^{\prime}\) 为函子。为使 \(\varphi\) 是同构，必须且只需映射 \(\operatorname{Som}(\varphi)\) 与 \(\operatorname{Fl}(\varphi)\) 都是双射。

这样，若把集合 S 与 F 上的范畴结构称为这些集合上的箭图结构与一个在每个顶点处存在单位元的结合合成律的数据，就可以把函子取作这个范畴结构的态射（E, IV, 第 11 页）。

